% Generated draft: requires mathematical and source review.
\documentclass[11pt]{article}
\usepackage[margin=25mm]{geometry}
\usepackage{fontspec}
\setmainfont{Latin Modern Roman}
\usepackage{amsmath,amssymb,mathtools,unicode-math}
\setmathfont{Latin Modern Math}
\usepackage{enumitem,needspace,xurl,hyperref}
\setlength{\parindent}{0pt}
\setlength{\parskip}{5pt}
\setlength{\emergencystretch}{3em}
\newcommand{\N}{\mathbb{N}}
\newcommand{\Z}{\mathbb{Z}}
\newcommand{\Q}{\mathbb{Q}}
\newcommand{\R}{\mathbb{R}}
\newcommand{\C}{\mathbb{C}}
\newcommand{\F}{\mathbb{F}}
\newcommand{\A}{\mathbb{A}}
\newcommand{\PP}{\mathbb{P}}
\newcommand{\E}{\mathbb{E}}
\newcommand{\eps}{\varepsilon}
\newcommand{\dd}{\,\mathrm{d}}
\begin{document}
\section*{UnsolvedMath ID: 800603}
\textbf{Source problem code:} catalogue entry
\subsection{Definitions and mathematical statement}
For a topological space \(Z\), an open cover is a family of open subsets whose union is \(Z\). Under the convention used below, the whole space \(Z\) is excluded as a member of a cover. Let \(\mathcal{O}(Z)\) denote these open covers. A cover \(\mathcal{U}\) is \emph{large} if every point of \(Z\) belongs to infinitely many distinct members of \(\mathcal{U}\); write \(\Lambda(Z)\) for the large open covers. The supplied record does not specify whether covers must be countable. These definitions use unrestricted families; imposing countability gives a potentially different formulation, and the intended convention requires source review.

The property \(\mathsf{Split}(\Lambda,\Lambda)\) means that every \(\mathcal{U}\in\Lambda(Z)\) admits a partition
\[
\mathcal{U}=\mathcal{U}_0\mathbin{\dot\cup}\mathcal{U}_1,
\qquad \mathcal{U}_0,\mathcal{U}_1\in\Lambda(Z).
\]
Thus, for each \(z\in Z\) and each \(i\in\{0,1\}\), infinitely many members of \(\mathcal{U}_i\) contain \(z\). A space is compact if every open cover has a finite subcover, and is \(\sigma\)-compact if it is a countable union of compact subspaces. Products carry the product topology.

\textbf{Question attributed in context to Zdomskyy (Problem 31).} If \(X\) is compact and \(Y\) satisfies \(\mathsf{Split}(\Lambda,\Lambda)\), must \(X\times Y\) satisfy \(\mathsf{Split}(\Lambda,\Lambda)\)? More explicitly, for the intended ambient class \(\mathcal{C}\) of spaces, the question is whether
\[
\forall X,Y\in\mathcal{C},\quad
\bigl[X\text{ is compact}\ \land\ Y\models\mathsf{Split}(\Lambda,\Lambda)\bigr]
\Longrightarrow
X\times Y\models\mathsf{Split}(\Lambda,\Lambda).
\]
The supplied record does not specify \(\mathcal{C}\), separation axioms, or restrictions such as metrizability. Consequently, this display is a formulation template with an unresolved ambient domain; it must not be read as evidence that the source asks the question for all topological spaces.

\textbf{Supplied contextual assertion.} To state the comparison included in the record, let \(\Omega(Z)\) consist of open covers such that every finite subset of \(Z\) is contained in some member, and let \(\Gamma(Z)\) consist of infinite open covers such that each point belongs to all but finitely many members. For \(\mathscr{B}\in\{\mathcal{O},\Omega,\Gamma\}\), the property \(\mathsf{U}_{\mathrm{fin}}(\mathcal{O},\mathscr{B})\) means that for every sequence \((\mathcal{U}_n)_{n\in\N}\) of open covers, none having a finite subcover, there are finite subfamilies \(\mathcal{F}_n\subseteq\mathcal{U}_n\) such that
\[
\left\{\bigcup\mathcal{F}_n:n\in\N\right\}\in\mathscr{B}(Z).
\]
The record states, with reference labels [SFT, AddQuad], that for each such \(\mathscr{B}\), the product of a \(\sigma\)-compact space \(X\) and a space \(Y\) satisfying \(\mathsf{U}_{\mathrm{fin}}(\mathcal{O},\mathscr{B})\) also satisfies that property. This is preserved as supplied background, subject to the unresolved ambient-space and cover conventions, and is not independently verified here.
\subsection{Short English statement}
Problem 31 asks whether taking a product with a compact space preserves the property that every large open cover can be partitioned into two large open covers. The record attributes the question in context to Zdomskyy. Its intended class of spaces and countability convention for covers remain unspecified. The registry labels the problem open; this is a supplied status, not an independently verified assessment.
\subsection{Sources}
\begin{itemize}
\item Supplied local UnsolvedMath registry record: permanent ID \(800603\), exact source problem code \texttt{catalogue entry}, in the AMR Open Problem Lists collection. Title: \emph{Selection Principles in Mathematics --- Problem 31 (Zdomskyy)}. The registry lists Tsaban as proposer and \(2006\) as the proposed year; its background gives the attribution/context as Zdomskyy.
\item Tsaban, \emph{Selection Principles and Special Sets of Reals: Open Problems} (2006), Problem 31. Supplied source URL: \url{https://arxiv.org/abs/math/0609601}. The record reports access on 2026-10-16 and extraction by \texttt{arxiv\_latex}; these are supplied provenance statements, not a claim of source access or verification in this draft.
\item The contextual product assertion carries the supplied reference labels [SFT, AddQuad]. Their bibliographic identities are not provided in the record and have not been reconstructed.
\item The supplied registry status is \texttt{open}, with record update dated 2026-10-16. This is supplied registry curation, not an independently verified literature assessment. No separate \texttt{literature\_checked\_at} field or \texttt{literature\_sources} links were supplied.
\end{itemize}
\end{document}
